% GENERATED FILE. Edit canonical lessons, then run npm run generate:books.
% canonical-source-hash: fnv1a64:56dd5134a07a1e12
% canonical-lessons: ZH-C377-jumin, ZH-C377-shequ, ZH-W377-qu-area, ZH-C377-xiaoqu, ZH-C377-zhouwei, ZH-C377-xuequfang

\chapter{Our Neighbourhood}
\label{ch:zh-our-neighbourhood}
\hlchaptermodality{377}

\begin{chapteropening}
\textbf{By the end of this chapter:} \emph{I can talk about the people who live here, our neighbourhood, a housing estate, what is around it and homes near good schools.}

Complete the last lesson of chapter 377: I can talk about the people who live here, our neighbourhood, a housing estate, what is around it and homes near good schools.
\end{chapteropening}

\section[jūmín]{jūmín — resident (of a place)}

\label{lesson:ZH-C377-jumin}



\begin{quote}
Before the new one: say the Mandarin for to give up one's seat, then the Mandarin for square dancing (in a public square).
\end{quote}

\subsection*{You'll want to know: jūmín}
\emph{jūmín} — \textquotedblleft{}resident (of a place)\textquotedblright{}.

\begin{quote}
\emph{zhèr de jūmín dōu hěn rèxīn} — The people who live here are all very kind.
\end{quote}

Its characters come later; for now it is a word for the ear.

\begin{cousinweb}
\emph{jūmín} is built from \emph{jū} 'to dwell' and \emph{mín} 'the people', both inherited from Middle Chinese, \emph{kjo} and \emph{mjin} (Guangyun, in Baxter's transcription), which glosses the first 'to stay' and the second 'the multitude'.
\end{cousinweb}

\begin{grammarlens}[title={What you've built}]
One more word for giving reasons.
\end{grammarlens}

\subsection*{Your turn}
\begin{itemize}
\raggedright
  \item \emph{Say it:} \emph{jūmín}
  \item \emph{Say it:} \emph{jūmín}, once more
  \item \emph{Say it:} \emph{zhèr de jūmín dōu hěn rèxīn} — The people who live here are all very kind.
  \item \emph{Recall:} say the Mandarin for to give up one's seat, then the Mandarin for square dancing (in a public square), then say \emph{jūmín} again
\end{itemize}

\begin{tcolorbox}[breakable,colback=teal!4,colframe=teal!35!black,title={Before you move on}]
What does jūmín mean? (\textquotedblleft{}Resident (of a place)\textquotedblright{}.) Say it once more.
\end{tcolorbox}
\section[\texorpdfstring{\zh{社区} (shèqū)}{社区 (shèqū)}]{\zh{社区} (shèqū) — community; neighbourhood}

\label{lesson:ZH-C377-shequ}



\begin{quote}
Before the new one: say the Mandarin for square dancing (in a public square), then the Mandarin for a resident (of a place).
\end{quote}

\subsection*{You'll want to know: \zh{社区}}
\textbf{\zh{社区}} — \emph{shèqū} — \textquotedblleft{}community; neighbourhood\textquotedblright{}.

\begin{quote}
\textbf{\zh{我们社区里有一家书店}} \emph{wǒmen shèqū lǐ yǒu yì jiā shūdiàn} — There's a bookshop in our neighbourhood.
\end{quote}

The new character, in pieces: \zh{区} — four strokes: a bar across the top, two strokes crossing inside, and a frame down the left and along the foot, open on the right. The next lesson puts it in your hand.

\begin{cousinweb}
\textbf{\zh{社区}} (\emph{shèqū}) is built from \zh{社} 'the altar of the soil' and \zh{区} 'an area', both inherited from Middle Chinese, \emph{dzyaeX} and \emph{khju} (Guangyun, in Baxter's transcription). The dictionary's entry for the second names a marsh and quotes a note on bent shoots; the sense 'an area' is not given there.
\end{cousinweb}

\begin{grammarlens}[title={What you've built}]
One more word for giving reasons.
\end{grammarlens}

\subsection*{Your turn}
\begin{itemize}
\raggedright
  \item \emph{Say it:} \emph{shèqū}
  \item \emph{Say it:} \emph{shèqū}, once more
  \item \emph{Say it:} \emph{wǒmen shèqū lǐ yǒu yì jiā shūdiàn} — There's a bookshop in our neighbourhood.
  \item \emph{Recall:} say the Mandarin for square dancing (in a public square), then the Mandarin for a resident (of a place), then say \emph{shèqū} again
\end{itemize}

\begin{tcolorbox}[breakable,colback=teal!4,colframe=teal!35!black,title={Before you move on}]
What does \zh{社区} mean? (\textquotedblleft{}Community; neighbourhood\textquotedblright{}.) Say it once more.
\end{tcolorbox}
\section[\texorpdfstring{\zh{区} (qū)}{区 (qū)}]{\zh{区} — a bar, a cross, then the frame}

\label{lesson:ZH-W377-qu-area}



\begin{quote}
The word first, then the character inside it.

\begin{itemize}
\raggedright
  \item \emph{Recall:} say \emph{shèqū} — \textbf{R1}, one lesson back, before it decays
\end{itemize}

\begin{quote}
\zh{社} \zh{小}
\end{quote}

\emph{Read it:} each one aloud, with nothing said first — \emph{shè}, \emph{xiǎo}
\end{quote}

\subsection*{Script — a bar, a cross, then the frame}
\hlblockfigure{\detokenize{figures/ZH-W377-qu-area-filmstrip.pdf}}{How \zh{区} is written, stroke by stroke}

\begin{quote}
\zh{区}
\end{quote}

\textbf{Three steps: the top bar, the cross inside, the frame.}

\textbf{The top bar, one stroke,} left to right across the top.

\textbf{The cross inside, two strokes.} A stroke falling to the lower left, then a stroke falling to the lower right that crosses it.

\textbf{The frame, one stroke.} A vertical down the left side that turns right along the foot. The right side stays open.

1 + 2 + 1 = four.

\subsection*{Your turn}
\begin{itemize}
\raggedright
  \item \emph{Write it:} \zh{区}, all four, with the model in front of you
  \item \emph{Cover:} write it again from memory, and count the pieces, not the strokes
\end{itemize}

\begin{tcolorbox}[breakable,colback=teal!4,colframe=teal!35!black,title={Before you move on}]
How many strokes in \zh{区}? (\textbf{Four}.) Which stroke comes last? (\textbf{The frame, down the left side and along the foot}.) Which side stays open? (\textbf{The right}.)

Source: \href{https://raw.githubusercontent.com/chanind/hanzi-writer-data/68d10a4b21150cae5e1ebbd223eed289cf32d90c/data/\%E5\%8C\%BA.json}{Hanzi Writer Data \zh{区}, pinned stroke paths and medians}.
\end{tcolorbox}
\section[\texorpdfstring{\zh{小区} (xiǎoqū)}{小区 (xiǎoqū)}]{\zh{小区} (xiǎoqū) — housing estate; residential compound}

\label{lesson:ZH-C377-xiaoqu}



\begin{quote}
Before the new one: say the Mandarin for a resident (of a place), then the Mandarin for a local community.
\end{quote}

\subsection*{You'll want to know: \zh{小区}}
\textbf{\zh{小区}} — \emph{xiǎoqū} — \textquotedblleft{}housing estate; residential compound\textquotedblright{}.

\begin{quote}
\textbf{\zh{我觉得这个小区很好}} \emph{wǒ juéde zhège xiǎoqū hěn hǎo} — I think this estate is very nice.
\end{quote}

\begin{cousinweb}
\textbf{\zh{小区}} (\emph{xiǎoqū}) is built from \zh{小} 'small' and \zh{区} 'an area', both inherited from Middle Chinese, \emph{sjewX} and \emph{khju} (Guangyun, in Baxter's transcription), which glosses the first 'tiny'. A housing estate of many blocks behind one gate is a modern use.
\end{cousinweb}

\begin{grammarlens}[title={What you've built}]
One more word for giving reasons.
\end{grammarlens}

\subsection*{Your turn}
\begin{itemize}
\raggedright
  \item \emph{Say it:} \emph{xiǎoqū}
  \item \emph{Say it:} \emph{xiǎoqū}, once more
  \item \emph{Say it:} \emph{wǒ juéde zhège xiǎoqū hěn hǎo} — I think this estate is very nice.
  \item \emph{Recall:} say the Mandarin for a resident (of a place), then the Mandarin for a local community, then say \emph{xiǎoqū} again
\end{itemize}

\begin{tcolorbox}[breakable,colback=teal!4,colframe=teal!35!black,title={Before you move on}]
What does \zh{小区} mean? (\textquotedblleft{}Housing estate; residential compound\textquotedblright{}.) Say it once more.
\end{tcolorbox}
\section[zhōuwéi]{zhōuwéi — surroundings; the area around}

\label{lesson:ZH-C377-zhouwei}



\begin{quote}
Before the new one: say the Mandarin for a local community, then the Mandarin for a housing estate.
\end{quote}

\subsection*{You'll want to know: zhōuwéi}
\emph{zhōuwéi} — \textquotedblleft{}surroundings; the area around\textquotedblright{}.

\begin{quote}
\emph{xiǎoqū zhōuwéi yǒu fàndiàn yě yǒu chāoshì} — Round the estate there are restaurants, and a supermarket too.
\end{quote}

Its characters come later; for now it is a word for the ear.

\begin{cousinweb}
\emph{zhōuwéi} is built from \emph{zhōu} 'all round' and \emph{wéi} 'to surround', both inherited from Middle Chinese, \emph{tsyuw} and \emph{hjw+j} (Guangyun, in Baxter's transcription), which glosses the first 'to go all the way round' and the second 'to encircle'.
\end{cousinweb}

\begin{grammarlens}[title={What you've built}]
One more word for giving reasons.
\end{grammarlens}

\subsection*{Your turn}
\begin{itemize}
\raggedright
  \item \emph{Say it:} \emph{zhōuwéi}
  \item \emph{Say it:} \emph{zhōuwéi}, once more
  \item \emph{Say it:} \emph{xiǎoqū zhōuwéi yǒu fàndiàn yě yǒu chāoshì} — Round the estate there are restaurants, and a supermarket too.
  \item \emph{Recall:} say the Mandarin for a local community, then the Mandarin for a housing estate, then say \emph{zhōuwéi} again
\end{itemize}

\begin{tcolorbox}[breakable,colback=teal!4,colframe=teal!35!black,title={Before you move on}]
What does zhōuwéi mean? (\textquotedblleft{}Surroundings; the area around\textquotedblright{}.) Say it once more.
\end{tcolorbox}
\section[\texorpdfstring{\zh{学区房} (xuéqūfáng)}{学区房 (xuéqūfáng)}]{\zh{学区房} (xuéqūfáng) — a home near a good school}

\label{lesson:ZH-C377-xuequfang}



\begin{quote}
Before the new one: say the Mandarin for a housing estate, then the Mandarin for the area around (a place).
\end{quote}

\subsection*{You'll want to know: \zh{学区房}}
\textbf{\zh{学区房}} — \emph{xuéqūfáng} — \textquotedblleft{}a home near a good school\textquotedblright{}.

\begin{quote}
\textbf{\zh{大家都想买学区房因为学校好}} \emph{dàjiā dōu xiǎng mǎi xuéqūfáng yīnwèi xuéxiào hǎo} — Everyone wants a flat near a good school, because the school is good.
\end{quote}

\begin{cousinweb}
\textbf{\zh{学区房}} (\emph{xuéqūfáng}) is built from \zh{学} 'to study', \zh{区} 'an area' and \zh{房} 'a house', all inherited from Middle Chinese, \emph{haewk}, \emph{khju} and \emph{bjang} (Guangyun, in Baxter's transcription), which glosses the last 'a room'. A home inside a good school's catchment area is a recent use.
\end{cousinweb}

\begin{grammarlens}[title={What you've built}]
One more word for giving reasons.
\end{grammarlens}

\subsection*{Your turn}
\begin{itemize}
\raggedright
  \item \emph{Say it:} \emph{xuéqūfáng}
  \item \emph{Say it:} \emph{xuéqūfáng}, once more
  \item \emph{Say it:} \emph{dàjiā dōu xiǎng mǎi xuéqūfáng yīnwèi xuéxiào hǎo} — Everyone wants a flat near a good school, because the school is good.
  \item \emph{Recall:} say the Mandarin for a housing estate, then the Mandarin for the area around (a place), then say \emph{xuéqūfáng} again
\end{itemize}

\begin{tcolorbox}[breakable,colback=teal!4,colframe=teal!35!black,title={Before you move on}]
What does \zh{学区房} mean? (\textquotedblleft{}A home near a good school\textquotedblright{}.) Say it once more.
\end{tcolorbox}
